\section*{Capítulo \ref{ch:b1:det} --- Determinantes e Sistemas Lineares}

\begin{solution}{exo:b1:det:1}
$3 \times 2 - 1 \times 5 = 1$.

Segundo: $L_2 \leftarrow L_2 - L_1$, $L_3 \leftarrow L_3 - L_2$ (sobre as linhas originais) dão as linhas
$(1,2,3), (3,3,3), (3,3,3)$: duas linhas iguais, determinante $0$. (Sarrus confirma:
$45 + 84 + 96 - 105 - 48 -
72 = 0$.)

Terceiro: é um Vandermonde com $x = 1, 2, 3$ (o \cref{ex:b1:det:vandermonde}): $(2-1)(3-1)(3-2) = 2$.
\end{solution}

\begin{solution}{exo:b1:det:2}
O determinante é igual a (some todas as colunas à primeira e
fatore) $(\lambda + 2)$ vezes
\[
\begin{vmatrix}
1 & 1 & \lambda\\ 1 & \lambda & 1\\ 1 & 1 & 1
\end{vmatrix}
= -(\lambda - 1)^2
\]
(zere com $L_1 \leftarrow L_1 - L_3$, $L_2 \leftarrow L_2 - L_3$ e desenvolva), o que dá $\det = -(\lambda+2)(\lambda-1)^2$. \omterm{def:b1:vspaces:free}{Base} $\iff \det \neq 0 \iff \lambda \notin \{1, -2\}$.
\end{solution}

\begin{solution}{exo:b1:det:3}
Primeiro sistema: $\det = -7$; $x = \frac{1}{-7}\begin{vmatrix} 5 & 1\\ 4 & -2\end{vmatrix}
= \frac{-14}{-7} = 2$, $y = \frac{1}{-7}\begin{vmatrix} 2 & 5\\ 3 & 4\end{vmatrix}
= \frac{-7}{-7} = 1$. Verificação: $2(2) + 1 = 5$; $3(2) - 2 = 4$.

Segundo sistema: depois de $L_2 - L_1$ e $L_3 - 2L_1$, as linhas tornam-se
$(1,1,1)$, $(0,-2,0)$, $(0,-1,-3)$, logo
\[
\det A = \begin{vmatrix} 1&1&1\\ 1&-1&1\\ 2&1&-1\end{vmatrix}
= 1 \times \begin{vmatrix} -2 & 0\\ -1 & -3\end{vmatrix} = 6 .
\]
Cramer, substituindo as colunas por $(6,2,1)^{\mathsf T}$:
\[
x = \frac{6}{6} = 1, \qquad
y = \frac{12}{6} = 2, \qquad
z = \frac{18}{6} = 3
\]
(os numeradores calculados do mesmo modo). Verificação: $1 + 2 + 3 = 6$;
$1 -
2 + 3 = 2$; $2 + 2 - 3 = 1$.
\end{solution}

\begin{solution}{exo:b1:det:4}
Reduza a matriz ampliada: $L_2 \leftarrow L_2 - 2L_1$, $L_3
\leftarrow L_3 - L_1$:
\[
\begin{pmatrix}
1 & 2 & -1 & 1 & 1\\
0 & 0 & 3 & -3 & 3\\
0 & 0 & 3 & -3 & 3
\end{pmatrix}
\to
\begin{pmatrix}
1 & 2 & -1 & 1 & 1\\
0 & 0 & 1 & -1 & 1\\
0 & 0 & 0 & 0 & 0
\end{pmatrix}.
\]
Incógnitas pivô $x, z$; incógnitas \omterm{def:b1:vspaces:free}{livres} $y, t$. Substituição
regressiva: $z =
1 + t$, $x = 1 - 2y + z - t = 2 - 2y$. \omterm{def:b1:logic:sets}{Conjunto} de soluções:
\[
\{(2 - 2y,\; y,\; 1 + t,\; t) : y, t \in \R\}
= (2, 0, 1, 0) + \operatorname{Vect}\bigl((-2,1,0,0),\,
(0,0,1,1)\bigr),
\]
um plano afim (dimensão $2 = 4 - \operatorname{rk} 2$) de $\R^4$.
\end{solution}

\begin{solution}{exo:b1:det:5}
Indução; $n = 1$ é o produto vazio $= 1$. Para o passo, efetue
$C_k \leftarrow C_k - x_1 C_{k-1}$ para $k = n, n-1, \dots, 2$ (nesta ordem, de modo que cada operação use uma
coluna ainda não modificada). A primeira linha torna-se $(1, 0, \dots, 0)$; na
linha $i \geq 2$, a $k$-ésima entrada torna-se $x_i^{k-1} - x_1 x_i^{k-2} = x_i^{k-2}(x_i - x_1)$. Desenvolvendo ao
longo da primeira linha e fatorando $(x_i - x_1)$ em cada linha $i$:
\[
V(x_1, \dots, x_n)
= \prod_{i=2}^{n} (x_i - x_1)\cdot V(x_2, \dots, x_n),
\]
e a hipótese de indução completa o produto $\prod_{i<j}(x_j
- x_i)$.
\end{solution}

\begin{solution}{exo:b1:det:6}
Desenvolvendo $D_n$ ao longo da primeira linha: $D_n = 2 D_{n-1} -
1\cdot\begin{vmatrix} 1 & \ast\\ 0 & D_{n-2}\text{-block}
\end{vmatrix}$; o segundo
determinante, desenvolvido ao longo da sua primeira coluna, é
$D_{n-2}$. Logo $D_n = 2D_{n-1} - D_{n-2}$, isto é, $D_n -
D_{n-1} = D_{n-1} - D_{n-2}$: as diferenças são constantes,
iguais a $D_2 - D_1 = 1$. Logo $D_n = D_1 + (n - 1) = n + 1$. (Verificação: $D_2 = 3$, e o caso $3\times3$
é o \cref{ex:b1:det:cofactor}: $D_3 = 4$.)
\end{solution}

\begin{solution}{exo:b1:det:7}
$C_1 \leftarrow C_1 + C_2 + C_3$ torna a primeira coluna constante e igual a $(m+2)$;
fatore-a:
\[
\det = (m+2)\begin{vmatrix}
1 & 1 & m\\ 1 & m & 1\\ 1 & 1 & 1
\end{vmatrix}
\overset{L_1 - L_3,\ L_2 - L_3}{=}
(m+2)\begin{vmatrix}
0 & 0 & m-1\\ 0 & m-1 & 0\\ 1 & 1 & 1
\end{vmatrix}
= (m+2)\cdot\bigl(-(m-1)^2\bigr)
\]
(desenvolva ao longo da primeira coluna: a única entrada $1$ carrega
sinal $+$, e o determinante $2 \times 2$ restante é $0 \cdot 0 -
(m-1)(m-1) = -(m-1)^2$).

Discussão. $m \notin \{1, -2\}$: sistema de Cramer; pela simetria das equações,
$x = y = z$, e cada equação dá $(m + 2)x = 1$: solução única $\bigl(\frac{1}{m+2}, \frac{1}{m+2},
\frac{1}{m+2}\bigr)$. $m = 1$: as
três equações leem-se todas $x + y
+ z = 1$: as soluções formam o plano afim
$x + y + z = 1$. $m = -2$: somar as três equações dá $0 = 3$: \omterm{def:b1:logic:sets}{conjunto} de
soluções vazio.
\end{solution}

\begin{solution}{exo:b1:det:8}
($\Rightarrow$) Se $A^{-1}$ tem entradas inteiras: $\det A \cdot \det
A^{-1} = 1$, com ambos os
determinantes inteiros (somas de produtos de entradas): dois
inteiros de produto $1$ são ambos $\pm1$.

($\Leftarrow$) A fórmula dos cofatores $A^{-1} = \frac{1}{\det
A}\operatorname{Com}(A)^{\mathsf T}$ (conferida para $n \leq 3$ por
desenvolvimento direto, admitida em geral) tem $\operatorname{Com}(A)$ com entradas
inteiras (cada cofator é um determinante inteiro); dividir por $\det A = \pm 1$
mantém os inteiros.
\end{solution}

\begin{solution}{exo:b1:det:9}
Some todas as colunas à primeira: cada entrada da nova primeira
coluna é $a + nb$; fatore-a, de modo que a primeira coluna fique só
de uns. Então as \omterm{met:b1:matrices:gauss}{operações sobre linhas} $L_i \leftarrow L_i - L_1$ ($i \geq 2$) zeram toda
entrada abaixo do $1$ superior esquerdo e deixam $a$ na diagonal e
$0$ no resto dessas linhas: a matriz é triangular superior com
diagonal $(1, a, \dots, a)$. Logo
\[
\det(aI + bJ) = (a + nb)\, a^{\,n-1} ,
\]
não nulo se e somente se $a \neq 0$ e $a + nb \neq 0$: a condição do \cref{exo:b1:matrices:9}.
\end{solution}

\begin{solution}{exo:b1:det:10}
Operações por blocos (cada uma é uma composição das $n$ operações
escalares correspondentes, permitidas pelo \cref{thm:b1:det:props}~(1)):
\[
\begin{vmatrix} A & B\\ B & A\end{vmatrix}
\overset{C_1 \leftarrow C_1 + C_2}{=}
\begin{vmatrix} A + B & B\\ A + B & A\end{vmatrix}
\overset{L_2 \leftarrow L_2 - L_1}{=}
\begin{vmatrix} A + B & B\\ 0 & A - B\end{vmatrix}
= \det(A+B)\,\det(A-B),
\]
usando a regra triangular por blocos no último passo.
\end{solution}

\begin{solution}{exo:b1:det:11}
$C_1 \leftarrow C_1 + C_2 + C_3$ torna a primeira coluna constante e igual a $(a + b + c)$; fatore-a, e
depois $L_2 \leftarrow L_2 - L_1$, $L_3 \leftarrow L_3 - L_1$:
\[
\Delta = (a+b+c)\begin{vmatrix}
1 & b & c\\
0 & a - b & b - c\\
0 & c - b & a - c
\end{vmatrix}
= (a+b+c)\bigl[(a-b)(a-c) + (b-c)^2\bigr],
\]
e, desenvolvendo, $(a-b)(a-c) + (b-c)^2 = a^2 + b^2 + c^2 - ab - bc
- ca$. Sobre $\C$, com $j^3 = 1$ e $1 + j + j^2 = 0$:
\begin{align*}
(a + jb + j^2c)(a + j^2b + jc)
&= a^2 + b^2 + c^2 + (j + j^2)(ab + bc + ca)\\
&= a^2 + b^2 + c^2 - ab - bc - ca ,
\end{align*}
donde a fatoração completa. (Estruturalmente: a coluna $(1, j, j^2)^{\mathsf T}$
satisfaz $M\,(1, j, j^2)^{\mathsf T}
= (a + jb + j^2c)(1, j, j^2)^{\mathsf T}$, e do mesmo modo para $j^2$ e $1$: os três fatores
são os três ``autovalores'' da circulante, uma história
sistematizada no volume do segundo ano de graduação.)
\end{solution}

\begin{solution}{exo:b1:det:12}
Escreva $r = \operatorname{rk} A$.

\emph{Existe uma submatriz $r \times r$ invertível.} Escolha $r$ colunas
\omterm{def:b1:vspaces:free}{livres} de $A$ e seja $B \in \mathcal{M}_{n,r}$ a matriz que elas formam: $\operatorname{rk} B = r$. Como o
posto por linhas é igual ao posto por colunas (o \cref{thm:b1:matrices:rank}), $B$ tem $r$
linhas \omterm{def:b1:vspaces:free}{livres}; conservar essas linhas dá uma submatriz $r \times r$ de $A$
de posto $r$, isto é, invertível.

\emph{Nenhuma maior existe.} Seja $S$ uma submatriz $s \times s$
invertível, tomada das colunas $j_1, \dots, j_s$ e das linhas $i_1,
\dots, i_s$ de $A$. Se
uma combinação $\sum_k \lambda_k
C_{j_k} = 0$ das colunas \emph{inteiras} correspondentes se
anula, então ler apenas as linhas $i_1, \dots, i_s$ dá $\sum_k
\lambda_k S_k = 0$ nas colunas de $S$,
logo todos os $\lambda_k = 0$ ($S$ invertível): as colunas $C_{j_1}, \dots, C_{j_s}$ de $A$ são
\omterm{def:b1:vspaces:free}{livres}, e $s \leq \operatorname{rk} A = r$.

Logo $\operatorname{rk} A$ é exatamente o maior tamanho de uma submatriz invertível.
\end{solution}

\begin{solution}{pb:b1:det:1}
\textbf{1.} $V(1,2,3,4) = (2-1)(3-1)(4-1)(3-2)(4-2)(4-3) = 1
\cdot 2\cdot 3\cdot 1\cdot 2\cdot 1 = 12$. A interpolação em nós distintos pede os
coeficientes de $P$ que resolvem $W c = y$ com $W = (x_i^{\,j-1})$, e $\det W = V \neq 0$: um
sistema de Cramer.

\textbf{2.} Percorra as colunas da esquerda para a direita. $C_1$ é
a coluna constante $P_0(x_i) = 1$ ($P_0$ \omterm{def:b1:poly:def}{mônico} de grau $0$). Suponha que as
colunas $1, \dots, j-1$ já foram reduzidas às potências puras $1, x_i, \dots, x_i^{\,j-2}$. Como
$P_{j-1}
= X^{j-1} + \sum_{k < j-1}\alpha_k X^k$, subtrair de $C_j$ a combinação $\sum_k \alpha_k\,(\text{coluna de } x_i^k)$ --- operação que não
altera o determinante --- deixa a coluna de potências puras $x_i^{\,j-1}$.
Depois da última coluna, a matriz é a matriz de Vandermonde: $\det = V(x_1, \dots, x_n)$.

\textbf{3.} Os \omterm{def:b1:poly:def}{polinômios} $(j-1)!\,B_{j-1}$ são \omterm{def:b1:poly:def}{mônicos} de grau $j - 1$, logo a
questão 2 dá
\[
\det\bigl(B_{j-1}(m_i)\bigr)
= \frac{V(m_1, \dots, m_n)}{0!\,1!\cdots(n-1)!} .
\]
O lado esquerdo é o determinante de uma matriz com entradas
\emph{inteiras} ($B_k$ assume valores inteiros em $\Z$: questões
16--17 do problema de fim de semana \cref{pb:b1:vspaces:1}), logo um inteiro; e ele é
positivo, pois $V(m_1, \dots, m_n) > 0$ para $m_1 <
\dots < m_n$. Logo o superfatorial \omterm{def:b1:arith:divides}{divide} o
produto de todas as diferenças dois a dois.

\textbf{4.} Fatore $x_i$ em cada linha $i$: $\det(x_i^{\,j})_{j = 1..n} = x_1\cdots x_n\,
\det(x_i^{\,j-1}) = x_1\cdots x_n\,V$.

\textbf{5.} $(W^{\mathsf T}W)_{ij} = \sum_k x_k^{\,i-1}
x_k^{\,j-1} = p_{i+j-2}$: $S = W^{\mathsf T}W$. Logo $\det S =
\det(W^{\mathsf T})\det W = V^2$ (o \cref{thm:b1:det:props}~(2),(4)). Para $x_i$
reais: $\det S = V^2
\geq 0$, e $S$ é invertível se e somente se $V \neq 0$, isto é, se e
somente se os $x_i$ são dois a dois distintos --- um teste de sinal
definido, computável apenas a partir das somas de potências.

\textbf{6.} As condições de interpolação $\sum_{k}
c_k\,x_i^{\,k} = y_i$ formam o sistema
$Wc = y$; $\det W = V \neq
0$ dá existência e unicidade de uma só vez. Esta é a
terceira demonstração do livro: fórmula explícita no \cref{thm:b1:poly:lagrange},
argumento de núcleo no \cref{ex:b1:linmaps:interpolation}, Cramer aqui.

\textbf{7.} Cramer: $c_{n-1} = \det W'/\det W$, onde $W'$ é $W$ com a sua última coluna
substituída por $y$. Desenvolvendo $\det W'$ ao longo dessa coluna:
\[
\det W' = \sum_{i=1}^n (-1)^{i+n} y_i\,V(x_1, \dots, \widehat{x_i},
\dots, x_n) .
\]
Ora, $V = V(\setminus i)\cdot\prod_{j<i}(x_i - x_j)\prod_{j>i}
(x_j - x_i)$, e converter o segundo produto custa $(-1)^{n-i}$:
\[
(-1)^{i+n}\,\frac{V(\setminus i)}{V}
= \frac{(-1)^{i+n}(-1)^{n-i}}{\prod_{j\neq i}(x_i - x_j)}
= \frac{1}{\prod_{j\neq i}(x_i - x_j)} ,
\]
donde $c_{n-1} = \sum_i y_i/\prod_{j \neq i}(x_i - x_j)$ --- de novo a fórmula das \omterm{pb:b1:vspaces:1}{diferenças divididas}.

\textbf{8.} $L_3 \leftarrow L_3 - L_1$ dá as linhas $(1, x_1,
x_1^2)$, $(0, 1, 2x_1)$, $(0,\ x_2 - x_1,\ (x_2-x_1)(x_2+x_1))$; desenvolvendo ao
longo da primeira coluna e fatorando $(x_2 - x_1)$:
\[
(x_2 - x_1)\begin{vmatrix} 1 & 2x_1\\ 1 & x_2 + x_1
\end{vmatrix} = (x_2 - x_1)(x_2 - x_1) = (x_2 - x_1)^2 .
\]
Não nulo para $x_1 \neq x_2$: o \omterm{def:b1:det:system}{sistema linear} que exprime $P(x_1) = u$, $P'(x_1) = v$,
$P(x_2) = w$ nos coeficientes de $P \in \R_2[X]$ é de Cramer --- a interpolação de
Hermite com um nó duplicado é bem posta.

\textbf{9.} $P = a + bX + cX^2$ com $a = P(0) = 1$, $b = P'(0)
= 0$, $a + b + c = P(1) = 2$: $c = 1$, logo $P = 1 + X^2$,
único. Coerência: aqui $x_1 = 0$, $x_2 = 1$ e o determinante da
questão 8 é $(1 - 0)^2 = 1 \neq 0$.

\textbf{10.} Cálculo direto:
\[
C_2 = \frac{1}{(a_1+b_1)(a_2+b_2)} - \frac{1}{(a_1+b_2)(a_2+b_1)}
= \frac{(a_1+b_2)(a_2+b_1) - (a_1+b_1)(a_2+b_2)}
{\prod_{i,j}(a_i+b_j)} ,
\]
e o numerador desenvolve-se em $a_1b_1 + a_2b_2 - a_1b_2 - a_2b_1
= (a_2 - a_1)(b_2 - b_1)$: diferenças sobre somas.

\textbf{11.} Para $i < n$, a nova entrada da linha $i$ é
\[
\frac{1}{a_i + b_j} - \frac{1}{a_n + b_j}
= \frac{a_n - a_i}{(a_i + b_j)(a_n + b_j)} .
\]
Fatore $(a_n - a_i)$ em cada linha $i < n$, e depois $\frac1{a_n + b_j}$ em cada coluna
$j$: o que resta tem entradas $\frac1{a_i + b_j}$ nas linhas $i < n$ e a constante
$1$ na linha $n$ --- a matriz $M$, com o prefator anunciado.

\textbf{12.} Em $M$, para $j < n$ a operação $C_j \leftarrow
C_j - C_n$ transforma a
linha $n$ em $(0, \dots, 0, 1)$ e, na linha $i < n$,
\[
\frac{1}{a_i + b_j} - \frac{1}{a_i + b_n}
= \frac{b_n - b_j}{(a_i + b_j)(a_i + b_n)} .
\]
Fatore $(b_n - b_j)$ em cada coluna $j < n$ e $\frac1{a_i +
b_n}$ em cada linha $i < n$, e
depois desenvolva ao longo da última linha (sinal $(-1)^{n+n} = +1$): o
determinante restante é $C_{n-1}$. Reunindo os fatores das questões
11--12:
\[
C_n = \frac{\prod_{i<n}(a_n - a_i)\,\prod_{j<n}(b_n - b_j)}
{\prod_{j}(a_n + b_j)\,\prod_{i<n}(a_i + b_n)}\;C_{n-1},
\]
e a indução (\omterm{def:b1:vspaces:free}{base} $C_1 = \frac1{a_1+b_1}$) monta exatamente o duplo alternante de
Cauchy: os fatores $(a_j - a_i)(b_j
- b_i)$ para todos os pares, sobre todas as somas
$(a_i + b_j)$.

\textbf{13.} A fórmula se anula se e somente se algum $a_j = a_i$ ou
$b_j =
b_i$: a matriz de Cauchy é invertível se e somente se as duas
famílias são dois a dois distintas. Hilbert: $a_i = i$, $b_j = j - 1$. Para
$n = 2$: numerador $(2-1)(1-0) = 1$, denominador $1\cdot2\cdot2\cdot3
= 12$: $\det H_2 = \frac1{12}$. Para $n = 3$:
numerador $\bigl[(1)(2)(1)\bigr]^2 = 4$, denominador $(1\cdot2\cdot3)
(2\cdot3\cdot4)(3\cdot4\cdot5) = 6\cdot24\cdot60 = 8640$: $\det H_3 = \frac{4}{8640} = \frac1{2160}$. Inversa para $n = 2$:
\[
H_2^{-1} = 12\begin{pmatrix} \frac13 & -\frac12\\[2pt]
-\frac12 & 1\end{pmatrix}
= \begin{pmatrix} 4 & -6\\ -6 & 12 \end{pmatrix},
\]
todos inteiros (fenômeno verdadeiro para todo $n$).

\textbf{14.} A matriz do sistema é a matriz de Cauchy, invertível
pela questão 13 quando os $b_j$ (e os $a_i$) são dois a dois
distintos: solução única. Interpretação: uma função racional $R = \sum_j \frac{c_j}{X + b_j}$
com \omterm{def:b1:fractions:field}{polos} simples é determinada por $n$ dos seus valores $R(a_1), \dots, R(a_n)$, e
reciprocamente qualquer tal folha de dados é realizada exatamente
uma vez --- a contrapartida amostral do teorema de existência e
unicidade das frações parciais (o \cref{thm:b1:fractions:complex}).

\textbf{15.} Viète para $X^3 + pX + q$: $\lambda_1 + \lambda_2 +
\lambda_3 = 0$, $\sum_{i<j}\lambda_i\lambda_j = p$, logo $p_1 = 0$ e $p_2 = p_1^2 - 2p = -2p$.
Cada raiz satisfaz $\lambda^3 =
-p\lambda - q$; somando: $p_3 = -p\,p_1 - 3q = -3q$. Multiplicando por
$\lambda$ e somando: $p_4 = -p\,p_2 - q\,p_1 = 2p^2$.

\textbf{16.} Pela questão 5 (a identidade $S = W^{\mathsf T}W$ e $\det S = V^2$ valem sobre
$\C$),
\[
V^2 = \begin{vmatrix}
3 & 0 & -2p\\
0 & -2p & -3q\\
-2p & -3q & 2p^2
\end{vmatrix}
= 3\bigl(-4p^3 - 9q^2\bigr) + (-2p)\bigl(0 - 4p^2\bigr)
= -4p^3 - 27q^2 ,
\]
desenvolvendo ao longo da primeira linha.

\textbf{17.} Uma raiz múltipla significa dois $\lambda_i$ iguais, isto é,
$V = 0$, isto é, $\operatorname{disc} = -4p^3 - 27q^2 = 0$. Para $X^3 - 3X + 2$: $4(-3)^3 + 27\cdot4 = -108 + 108 = 0$, coerente com a raiz dupla
$1$ de $(X-1)^2(X+2)$.

\textbf{18.} As raízes não reais de uma cúbica real vêm em pares
conjugados, logo ocorrem exatamente dois casos quando $\operatorname{disc} \neq
0$. Três
raízes reais distintas: $V$ é real e não nulo, logo $\operatorname{disc} = V^2 > 0$. Uma raiz
real $\lambda_1$ e $\lambda_3 = \conj{\lambda_2} \notin \R$: então
\[
(\lambda_2 - \lambda_1)(\lambda_3 - \lambda_1) =
\abs{\lambda_2 - \lambda_1}^2 > 0,
\qquad
\lambda_3 - \lambda_2 = -2\iu\,\operatorname{Im}\lambda_2 \neq 0,
\]
logo $V$ é um número imaginário puro não nulo e $\operatorname{disc} = V^2 < 0$. Os dois
sinais caracterizam os dois casos.

\textbf{19.} Tome $b_i = a_i$ no duplo alternante: o numerador é $\prod_{i<j}(a_j - a_i)^2 > 0$ e o
denominador $\prod_{i,j}(a_i + a_j) > 0$ (todas as entradas positivas): o determinante é
positivo. (Em linguagem posterior: o núcleo $\frac1{x+y}$ é positivo
definido.)

\textbf{20.} Pelas questões 2--3, o determinante vale $V(2,4,7)/(0!\,1!\,2!) = \frac{(4-2)(7-2)(7-4)}{2} =
\frac{30}{2} = 15$.
Diretamente, a matriz é
\[
\begin{pmatrix}
1 & 2 & 1\\
1 & 4 & 6\\
1 & 7 & 21
\end{pmatrix},
\qquad
\det = (84 - 42) - 2(21 - 6) + (7 - 4) = 42 - 30 + 3 = 15 .
\]

\textbf{21.} Suponha $\sum_i c_i\,(\lambda_i^{\,k})_{k} = 0$ como sequência. Ler em $k = 0, 1, \dots, n-1$ dá $W^{\mathsf
T}c = 0$ com
$W = (\lambda_i^{\,j-1})$ invertível ($\det = V
\neq 0$, com $\lambda_i$ distintos): $c = 0$. As
sequências geométricas são \omterm{def:b1:vspaces:free}{livres}.

\textbf{22.} $a = b = (1, 2, 3)$: numerador $\bigl[(2-1)(3-1)
(3-2)\bigr]^2 = 4$; denominador $\prod_{i,j}(i + j) =
(2\cdot3\cdot4)(3\cdot4\cdot5)(4\cdot5\cdot6) = 24\cdot60\cdot120
= 172800$. Logo $\det\bigl(\frac1{i+j}\bigr) = \frac{4}{172800}
= \frac1{43200}$.

\textbf{23.} Se $x_i = x_j$, a troca das duas variáveis fixa o ponto mas
deve mudar o sinal de $F$: $F = -F$, logo $F = 0$ ali.
\omterm{def:b1:arith:divides}{Divisibilidade}: veja $F$ como um \omterm{def:b1:poly:def}{polinômio} na única variável $x_n$
com coeficientes nas outras variáveis; ele se anula nos $n - 1$
``valores'' $x_1, \dots, x_{n-1}$, logo fatorações repetidas (o \cref{thm:b1:poly:factor}) dão $F =
\prod_{i<n}(x_n - x_i)\cdot G$ com
$G$ polinomial. O prefator é invariante pelas trocas de dois índices
$i, j < n$, logo $G$ é alternante em $x_1, \dots, x_{n-1}$, e a indução completa: $\prod_{i<j}(x_j - x_i)$
\omterm{def:b1:arith:divides}{divide} $F$.

\textbf{24.} $D = \det(x_i^{\,j-1})$ é um \omterm{def:b1:poly:def}{polinômio} nos $x_i$; trocar duas variáveis
troca duas linhas, logo $D$ é alternante e, pela questão 23, $D = c\,\prod_{i<j}(x_j - x_i)$
para algum \omterm{def:b1:poly:def}{polinômio} $c$. Graus totais: $D$ tem grau $\leq 0 +
1 + \dots + (n-1) = \binom n2$, e o
produto tem grau exatamente $\binom n2$: $c$ é uma constante. O monômio
$x_2\,x_3^2\cdots
x_n^{\,n-1}$ tem coeficiente $1$ em $D$ (produto diagonal) e $1$ no produto
(escolha a variável de índice maior em cada fator): $c = 1$, e a
fórmula de Vandermonde sai sem indução alguma.

\textbf{25.} (i) A alternância --- o axioma ``duas colunas iguais
matam o determinante'' --- é o motor: foi ela que produziu cada
fator $(x_j - x_i)$, $(a_j - a_i)$, $(b_j - b_i)$ do problema. (ii) A identidade $\det S = V^2$
substitui as raízes complexas individuais, inatingíveis, pelas suas
somas de potências, que são \omterm{def:b1:poly:def}{polinômios} reais nos coeficientes, de
modo que a distinção se torna o sinal de um número real computável.
(iii) O \omterm{ex:b1:det:vandermonde}{determinante de Vandermonde} governa a interpolação
polinomial, e o \omterm{pb:b1:det:1}{determinante de Cauchy} governa as frações parciais e
as funções racionais amostradas (com a \omterm{pb:b1:det:1}{matriz de Hilbert} como o seu
caso particular mais famoso). (iv) Todo \omterm{def:b1:poly:def}{polinômio} alternante é
divisível por $\prod_{i<j}(x_j - x_i)$, e uma contagem de graus fixa então tal
\omterm{def:b1:poly:def}{polinômio} a menos de uma constante --- razão pela qual este produto
não para de reaparecer sempre que um determinante se anula sobre
coincidências. O teorema da Parte III é o \emph{duplo alternante de
Cauchy}.
\end{solution}
